\documentclass[12pt,a4paper,fontset=none]{ctexart}
\usepackage[margin=1in]{geometry}
\usepackage{setspace}
\usepackage{booktabs}
\usepackage{longtable}
\usepackage{array}
\usepackage{amsmath,amssymb}
\usepackage{hyperref}
\usepackage{enumitem}

% English text: prefer Times New Roman, fall back to a Times-compatible font.
\IfFontExistsTF{Times New Roman}
  {\setmainfont{Times New Roman}}
  {\setmainfont{TeX Gyre Termes}}

% Use standard macOS font for Chinese text
\setCJKmainfont[AutoFakeBold=2,AutoFakeSlant=.2]{PingFang TC}
\setCJKsansfont[AutoFakeBold=2,AutoFakeSlant=.2]{PingFang TC}
\setCJKmonofont{PingFang TC}
\renewcommand{\contentsname}{目錄}

\setstretch{1.2}
\setlength{\parskip}{0.5em}
\setlength{\parindent}{2em}
\newcolumntype{L}[1]{>{\raggedright\arraybackslash}p{#1}}

\title{How Textual Features Influence Social Media Engagement:\\Evidence from Booking.com on X\\[0.5em]
\large Booking.com 社群文字特徵與互動成效之研究報告（雙語版）\\[0.5em]
\small v1.2.0 - Descriptive Statistics \& Correlation Tables Added}
\author{}
\date{\today}

\begin{document}
\maketitle
\tableofcontents
\newpage

\section*{Part I. English Version}
\addcontentsline{toc}{section}{Part I. English Version}

\section{Background and Motivation}
Social media engagement has become an important performance indicator for brands because likes, comments, and shares reflect audience attention, interaction, and the social diffusion of branded content. For a travel platform such as Booking.com, engagement on X can be especially valuable because travel decisions are socially visible, information-intensive, and often shaped by peer endorsement. Mohammad et al. (2024) argue that social commerce interactions can strengthen credibility and influence downstream behavioral outcomes, which makes engagement a meaningful intermediate outcome for hospitality and travel brands. In this context, understanding which textual features encourage stronger engagement can help Booking.com design more effective posts and communicate more persuasively with potential travelers.

This proposal focuses on the textual design of Booking.com's posts rather than on broad brand strategy. The central idea is that the way a post is written can affect how users react to it. Prior literature suggests that content structure, interactivity cues, emotional tone, and discoverability signals can influence user responses on social media (Moran et al., 2019; van der Harst and Angelopoulos, 2024). Building on that logic, this study examines whether text length, question marks, emotional valence, and hashtags are associated with higher engagement for Booking.com posts on X, while controlling for non-textual post elements such as pictures (Li and Xie, 2020), URLs, and at-mentions.

\section{Research Question}
The research question guiding this study is:

\begin{quote}
How do textual features in Booking.com's posts on X influence social media engagement?
\end{quote}

More specifically, the study asks whether longer posts, more interactive wording through question marks, more positive emotional valence, and a greater number of hashtags are associated with higher engagement, after accounting for other post characteristics. This question is theoretically relevant because it links content design decisions to observable consumer response, and it is managerially relevant because it can guide the brand's day-to-day social media writing strategy.

\section{Hypotheses}
The hypotheses are derived from prior studies on social media engagement and content design. First, longer text may provide richer information, context, and persuasive detail, which can help users understand the message and increase their willingness to engage. Gkikas et al. (2022) show that text characteristics, including length, are relevant predictors of engagement in social media contexts. In a travel setting, longer copy may also better communicate destinations, experiences, and value propositions.

Second, question marks are treated as an interactivity cue. De Vries, Gensler, and Leeflang (2012) argue that questions invite audience participation because they implicitly ask users to respond. For a brand such as Booking.com, a question can encourage users to imagine travel choices, express preferences, or reply with opinions, which should raise engagement.

Third, emotional valence captures whether the wording of a post is more positive or negative. Reddy and Prakash (2024) suggest that more positive sentiment can enhance user interaction and loyalty in social media environments. Because tourism communication often relies on aspiration, excitement, and pleasant imagery, positively valenced language is expected to generate stronger reactions.

Fourth, hashtags may improve discoverability, organize topical meaning, and connect a post to broader conversations. Gkikas et al. (2022) demonstrate that text characteristics, including the number of hashtags, positively impact user engagement rates outperforming other post features. For Booking.com, hashtags may make travel posts easier to locate and more visible within trend-based conversations, thereby improving engagement.

Fifth, emoji characters can function as visual-linguistic cues that convey emotion, add personality, and increase the visual salience of a post. Ko et al. (2022) show that emoji usage in brand-related content positively influences user engagement by making posts more expressive and approachable. In a travel context, emoji may complement textual descriptions of destinations and experiences, drawing reader attention and encouraging interaction.

Based on this reasoning, the study proposes the following hypotheses:
\begin{enumerate}[label=H\arabic*.]
    \item Text length is positively associated with social media engagement.
    \item Question marks are positively associated with social media engagement.
    \item Emotional valence is positively associated with social media engagement.
    \item The number of hashtags is positively associated with social media engagement.
    \item Emoji usage is positively associated with social media engagement.
\end{enumerate}

\section{Data}
The study uses the cleaned dataset in \texttt{data/Final\_Report\_data\_cleaned (2) (1).xlsx}. The dataset contains 766 posts from Booking.com's official X account. Based on the timestamp field, the observed period runs from November 1, 2021 to November 21, 2022. The dataset includes the original post text, engagement counts, media-related indicators, and multiple preprocessed text fields generated for text analysis.

The dependent variable is a combined engagement measure defined as the sum of likes, comments, and shares. This definition is consistent with the group's current research design and is grounded in the idea that engagement reflects multiple forms of observable user interaction. In the present sample, engagement is highly right-skewed: the median combined engagement is 4, the maximum is 2,474, and 104 posts receive zero combined engagement. This pattern is typical for social media data and has implications for model choice.

The dataset is suitable for the proposed analysis because it already includes the main textual and control variables required for the study, including length, question marks, hashtags, emotional valence, pictures, URLs, at-mentions, and emoji counts.

\section{Measurement of Variables}
This study operationalizes all variables using fields that already exist in the cleaned dataset and the preprocessing logic documented in \texttt{py\_export/single/final\_report\_full.py}. The main variables are defined as follows.

\small
\begin{longtable}{@{}L{0.20\textwidth} L{0.18\textwidth} L{0.36\textwidth} L{0.14\textwidth}@{}}
\toprule
Construct & Role & Operational definition & Dataset field \\
\midrule
\endfirsthead
\toprule
Construct & Role & Operational definition & Dataset field \\
\midrule
\endhead
Social media engagement & Dependent variable & Sum of likes, comments, and shares for each post & Derived from \texttt{like}, \texttt{comment}, \texttt{share} \\
Text length & Main independent variable & Number of cleaned English words after removing URLs, mentions, hashtags, and non-letter symbols, then filtering to recognized English words & \texttt{length} \\
Question mark & Main independent variable & Binary dummy tracking if the post contains at least one question mark after removing URLs (0=No, 1=Yes) & \texttt{question} \\
Emotional valence & Main independent variable & Lexicon-based sentiment score centered around neutral tone; higher values indicate more positive wording & \texttt{valence} \\
Hashtag & Main independent variable & Number of hashtags contained in the post & \texttt{hashtag} \\
Picture & Control variable & Number of pictures or visual media items attached to the post & \texttt{picture} \\
URL & Control variable & Number of URLs contained in the post & \texttt{url} \\
At-mention & Control variable & Number of \texttt{@account} mentions contained in the post & \texttt{at\_mention} \\
Emoji & Main independent variable & Number of emoji characters contained in the post & \texttt{emoji} \\
\bottomrule
\end{longtable}
\normalsize

The variable definitions are theoretically motivated and empirically grounded in the data preparation workflow. Text length captures how much verbal information a post provides. The \texttt{length} field is based on cleaned text rather than raw character count, which is appropriate because it better reflects the substantive amount of textual content available to readers. The question-mark variable captures a direct interactive cue and is operationalized as a binary dummy variable, since the vast majority of interrogative posts contain only a single question mark.

Emotional valence is measured with the existing \texttt{valence} field, which is generated through lexicon matching on the cleaned tokens and expresses the positivity or negativity of the language used in each post. This proposal treats valence as a continuous sentiment measure rather than splitting sentiment into separate positive and negative variables at this stage. Hashtags are measured as counts because the number of hashtags may reflect stronger discoverability or stronger topic signaling. Emoji is measured as a count of emoji characters in each post; prior literature suggests that emoji can increase engagement by adding expressive and visual-linguistic cues to branded content (Ko et al., 2022). The control variables capture non-textual or structural post elements that may affect engagement independently of writing style.

\section{Proposed Model Specification}
The empirical model examines how the five focal textual features relate to social media engagement while controlling for other post characteristics:

\begin{equation}
\begin{aligned}
\log(1 + \text{Engagement}_i) = {}& \beta_0 + \beta_1\text{TextLength}_i + \beta_2\text{Question}_i + \beta_3\text{Valence}_i \\
& + \beta_4\text{Hashtag}_i + \beta_5\text{Emoji}_i + \beta_6\text{Picture}_i + \beta_7\text{URL}_i \\
& + \beta_8\text{AtMention}_i + \varepsilon_i
\end{aligned}
\end{equation}

Conceptually, this specification tests whether the proposed textual features explain variation in engagement above and beyond visual or structural elements in the post. Because the dependent variable is a non-negative count-based sum and the sample shows substantial right skew, the dependent variable is modeled as
\[
\log\bigl(1 + \text{Engagement}_i\bigr)
\]
to reduce the influence of extreme observations and improve interpretability. The conceptual definition of engagement, however, remains the raw sum of likes, comments, and shares.

\section{Empirical Results}
Prior to testing the hypotheses, diagnostic checks were performed on the regression model. A Breusch-Pagan test revealed the presence of heteroscedasticity ($p < .0001$), while a Shapiro-Wilk test indicated that the residuals deviate from a perfectly normal distribution ($p < .0001$). However, given the large sample size ($N = 766$), ordinary least squares (OLS) regression relies on asymptotic normality. To correct for heteroscedasticity, we estimated the model using Heteroscedasticity-Consistent (HC) standard errors.

Table \ref{tab:desc} presents the descriptive statistics of all variables. Table \ref{tab:corr} presents the Pearson correlation matrix (lower triangle only). Table \ref{tab:results} presents the results of the robust OLS estimation.

The model explains 49.2\% of the variance in social media engagement. Contrary to Hypothesis 1, text length was found to have a significant \textit{negative} association with engagement ($\beta = -0.021, p < .001$), thus rejecting H1. As predicted, the presence of a question mark significantly increases engagement ($\beta = 0.479, p < .001$), supporting H2. Higher positive emotional valence ($\beta = 0.084, p < .01$) and a higher number of hashtags ($\beta = 1.096, p < .001$) both demonstrate significant positive relationships with engagement, supporting H3 and H4. Emoji usage also exhibits a statistically significant positive effect on engagement ($\beta = 0.343, p < .001$), supporting H5 and indicating that visual-linguistic cues enhance user interaction. 

\section{Discussion and Managerial Implications}
This study contributes to the literature on social media content strategy by evaluating textual design feature impacts within the context of a highly visual and fast-paced travel brand on X. By validating established content-feature theories, our study yields distinct managerial insights tailored precisely for travel platform marketing.

The rejection of H1 provides a highly intriguing revelation for Booking.com social media managers: on a micro-blogging platform like X, excessive text volume actively harms engagement. Travel audiences value succinct, quickly digestible information. In stark contrast, interactive elements significantly amplify engagement. Incorporating a question mark directly invites user participation and fosters a conversational dynamic. Similarly, portraying a positive emotional tone resonates strongly with the aspirational nature of tourism, and strategically deploying hashtags enhances thread discoverability amidst trending travel conversations. The support for H5 further confirms that emoji, as visual-linguistic cues, enrich branded content by conveying expressiveness and drawing reader attention, thereby boosting engagement.

\section{Conclusion}
In conclusion, the configuration of textual features dictates engagement outcomes on X. Our primary recommendation for Booking.com is to craft concise, overwhelmingly positive posts anchored closely with a question to immediately spur community interaction. Strategically adopting hashtags and injecting descriptive emojis further secures user attention without needlessly extending the text structure.

While our dependent variable amalgamates likes, comments, and retweets into a combined engagement metric, it is clear that optimizing textual interactivity operates as a foundational gateway for stimulating overall audience participation. Future research could segment engagement types or examine whether certain destination portrayals mediate these emotional text outcomes, but for immediate strategic alignment, the current findings offer a robust benchmark for daily social media copywriting.

\section{References (Selected)}
\begin{enumerate}
    \item De Vries, L., Gensler, S. and Leeflang, P.S.H. (2012). Popularity of brand posts on brand fan pages: An investigation of the effects of social media marketing. \textit{Journal of Interactive Marketing}, 26(2), 83--91.
    \item Gkikas, D. C., Tzafilkou, K., Theodoridis, P. K., Garmpis, A. and Gkikas, M. C. (2022). How do text characteristics impact user engagement in social media posts? Modeling content readability, length, and hashtags number in Facebook. \textit{International Journal of Information Management Data Insights}, 2, 100067.
    \item Ko, E.E., Kim, D. and Kim, G. (2022). Influence of emojis on user engagement in brand-related user-generated content. \textit{Computers in Human Behavior}.
    \item Li, Y. and Xie, Y. (2020). Is a picture worth a thousand words? An empirical study of image content and social media engagement. \textit{Journal of Marketing Research}, 57(1), 1--19.
    \item Mohammad, A.A., Elshaer, I.A., Azazz, A.M., Kooli, C., Algezawy, M. and Fayyad, S. (2024). The influence of social commerce dynamics on sustainable hotel brand image, customer engagement, and booking intentions. \textit{Sustainability}, 16(14), 6050.
    \item Moran, G., Muzellec, L. and Johnson, D. (2019). Message content features and social media engagement: evidence from the media industry. \textit{Journal of Product and Brand Management}, 29(5), 533--545.
    \item Reddy, M.A.S. and Prakash, C. (2024). Sentiment Analysis of Social Media Networking Sites: A Comparative Study on User Engagement and Public Opinion.
    \item van der Harst, J.P. and Angelopoulos, S. (2024). Less is more: Engagement with the content of social media influencers. \textit{Journal of Business Research}, 181, 114746.
\end{enumerate}

\newpage
\section*{Part II. 繁體中文版本}
\addcontentsline{toc}{section}{Part II. 繁體中文版本}

\section{研究背景與動機}
社群媒體互動（engagement）已成為品牌重要的績效指標，因為按讚、留言與分享能反映受眾注意力、互動程度，以及品牌內容在社群中的擴散效果。對 Booking.com 這類旅遊平台而言，X 上的互動更具價值，因為旅遊決策高度社會化、資訊需求高，且常受到同儕意見影響。Mohammad 等人（2024）指出，社群商務互動能提升品牌可信度，並影響後續行為結果，因此對旅宿與旅遊品牌而言，互動可視為具管理意義的中介成果。

本研究聚焦於 Booking.com 貼文的「文字設計」，而非整體品牌策略。核心概念是：貼文怎麼寫，會影響使用者如何反應。既有文獻顯示，內容結構、互動線索、情緒語調與可發現性訊號，都可能影響社群回應（Moran 等人，2019；van der Harst 與 Angelopoulos，2024）。據此，本研究檢驗文字長度、問號、情緒正向程度與 hashtag 是否與更高互動有關，並控制圖片（Li 與 Xie，2020）、URL 與 at-mention 等非文字因素。

\section{研究問題}
本研究的核心問題如下：

\begin{quote}
Booking.com 在 X 上的貼文文字特徵，如何影響社群互動表現？
\end{quote}

更具體地說，本研究探討：在控制其他貼文特徵後，較長文字、較多問號、較正向情緒語氣與較多 hashtag，是否與更高互動顯著相關。此問題在理論上連結了內容設計決策與可觀察的消費者反應；在管理上則可直接指引品牌日常社群文案策略。

\section{研究假說}
本研究假說建構自社群互動與內容設計相關文獻。首先，較長文字可能提供更豐富資訊、脈絡與說服細節，協助受眾理解訊息並提高互動意願。Gkikas 等人（2022）指出，文字特徵（含長度）是社群互動的重要預測因子。在旅遊情境中，較長文案也更有機會完整傳達目的地體驗與價值主張。

其次，問號可視為互動提示。De Vries、Gensler 與 Leeflang（2012）認為，提問會隱含邀請受眾回應，因而促進參與。對 Booking.com 而言，提問可鼓勵使用者想像旅遊選擇、表達偏好或留言回覆，進而提升互動。

第三，情緒正向程度（valence）反映貼文語句偏正面或偏負面。Reddy 與 Prakash（2024）指出，較正向情緒有助於提升社群互動與忠誠。在旅遊溝通中，由於訊息常依賴嚮往、興奮與愉悅想像，正向語氣預期能帶來更強反應。

第四，hashtag 可能提升可發現性、組織主題意義，並將貼文連結至更廣泛對話。Gkikas 等人（2022）實證指出，貼文特徵中的 hashtag 數量對提升社群互動率有顯著的正向貢獻。對 Booking.com 而言，hashtag 有助於旅遊內容被搜尋與被看見，從而提高互動。

第五，emoji 字元可作為視覺語言線索，傳達情感、增添個性，並提升貼文的視覺顯著度。Ko 等人（2022）指出，品牌相關內容中的 emoji 使用能正向影響互動，使貼文更具表現力與親近感。在旅遊情境中，emoji 可補充目的地與體驗的文字描述，吸引讀者目光並鼓勵互動。

基於上述推論，提出以下假說：
\begin{enumerate}[label=H\arabic*.]
    \item 文字長度與社群互動呈正向關係。
    \item 問號使用與社群互動呈正向關係。
    \item 情緒正向程度與社群互動呈正向關係。
    \item hashtag 數量與社群互動呈正向關係。
    \item emoji 使用與社群互動呈正向關係。
\end{enumerate}

\section{資料說明}
本研究使用清理後資料集 \texttt{data/Final\_Report\_data\_cleaned (2) (1).xlsx}，共包含 Booking.com 官方 X 帳號 766 筆貼文。依時間戳記欄位，觀察期間為 2021 年 11 月 1 日至 2022 年 11 月 21 日。資料含原始貼文文字、互動數據、媒體指標與多項文字前處理欄位。

依變數為「綜合互動量」，定義為 likes、comments、shares 三者加總。此定義符合目前小組研究設計，並反映互動是多種可觀察行為的集合。樣本中互動分布明顯右偏：中位數為 4、最大值為 2,474，且有 104 筆貼文的綜合互動為 0。此型態符合社群資料常見特徵，也影響模型選擇。

資料已包含本研究主要自變數與控制變數，包括 length、question、hashtag、valence、picture、url、at-mention、emoji，適合直接進行後續模型估計。

\section{變數衡量}
本研究所有變數皆以現有清理資料欄位操作化，並對應 \texttt{py\_export/single/final\_report\_full.py} 的前處理邏輯。主要變數定義如下。

\small
\begin{longtable}{@{}L{0.18\textwidth} L{0.16\textwidth} L{0.40\textwidth} L{0.14\textwidth}@{}}
\toprule
構念 & 角色 & 操作型定義 & 資料欄位 \\
\midrule
\endfirsthead
\toprule
構念 & 角色 & 操作型定義 & 資料欄位 \\
\midrule
\endhead
社群互動量 & 依變數 & 每篇貼文 likes、comments、shares 之加總 & 由 \texttt{like}、\texttt{comment}、\texttt{share} 推導 \\
文字長度 & 主要自變數 & 移除 URL、mentions、hashtags 與非字母符號後，再過濾為可辨識英文詞彙之詞數 & \texttt{length} \\
問號 & 主要自變數 & 虛擬變數（若貼文移除 URL 後含有一或多個問號則為 1，否則為 0） & \texttt{question} \\
情緒正向程度 & 主要自變數 & 詞典法情緒分數，以中性為中心；值越高代表語氣越正向 & \texttt{valence} \\
hashtag 數量 & 主要自變數 & 貼文中 hashtag 出現次數 & \texttt{hashtag} \\
圖片 & 控制變數 & 貼文附帶圖片或視覺媒體數量 & \texttt{picture} \\
URL & 控制變數 & 貼文中 URL 數量 & \texttt{url} \\
at-mention & 控制變數 & 貼文中 \texttt{@account} 提及次數 & \texttt{at\_mention} \\
emoji & 主要自變數 & 貼文中 emoji 字元數量 & \texttt{emoji} \\
\bottomrule
\end{longtable}
\normalsize

上述變數設定兼具理論與實務可行性。\texttt{length} 以清理後詞數而非原始字元數衡量，較能反映讀者可接收的實質文字資訊量。考量樣本中絕大多數含疑問語氣之貼文僅有單一問號，具有高度偏態，因此問號特徵以虛擬變數（dummy variable）處理，將重點聚焦於互動意圖的有無。\texttt{valence} 以連續分數呈現語言正負向程度，暫不拆分為正面與負面兩個獨立構面。hashtag 以計數表示主題標記與可發現性強度。emoji 以計數衡量貼文中表情符號的使用量；既有文獻指出 emoji 可作為視覺語言線索，增添品牌內容的表現力與親近感，進而提升互動（Ko 等人，2022）。控制變數用以捕捉文字以外的結構因素。

\section{模型設定}
本研究以下列模型檢驗五個核心文字特徵對互動的影響，並控制其他貼文特徵：

\begin{equation}
\begin{aligned}
\log(1 + \text{Engagement}_i) = {}& \beta_0 + \beta_1\text{TextLength}_i + \beta_2\text{Question}_i + \beta_3\text{Valence}_i \\
& + \beta_4\text{Hashtag}_i + \beta_5\text{Emoji}_i + \beta_6\text{Picture}_i + \beta_7\text{URL}_i \\
& + \beta_8\text{AtMention}_i + \varepsilon_i
\end{aligned}
\end{equation}

此模型用以檢驗：文字特徵是否能在視覺與結構因素之外，額外解釋互動差異。由於依變數為非負計數加總且右偏明顯，本研究以
\[
\log\bigl(1 + \text{Engagement}_i\bigr)
\]
作為實證估計之依變數，以降低極端值影響並提升解讀性；互動的概念定義仍維持 likes、comments、shares 的原始加總。

\section{實證結果}
在測試假說前，本研究首先針對迴歸模型進行了診斷檢驗。Breusch-Pagan 檢定結果顯示存在異質變異數 ($p < .0001$)，而 Shapiro-Wilk 檢定顯示殘差非常態分配 ($p < .0001$)。儘管如此，基於較大的樣本數 ($N = 766$)，一般最小平方估計 (OLS) 在漸近常態特性下依然成立。為校正異質變異數所可能引發的偏誤，模型採用「穩健標準誤」(Robust Standard Errors, HC) 進行估計。

表 \ref{tab:desc} 呈現所有變數的敘述統計。表 \ref{tab:corr} 呈現 Pearson 相關矩陣（僅顯示下三角）。表 \ref{tab:results} 呈現了穩健 OLS 估計結果。

模型所解釋的社群互動變異量為 49.2\%。與假說一的預期相反，文字長度對社群互動有顯著的「負向」關聯 ($\beta = -0.021, p < .001$)，因此拒絕 H1。如同假說預期，貼文中包含問號標點符號會顯著提振互動 ($\beta = 0.479, p < .001$)，支持 H2。此外，越強烈的正向情緒用語 ($\beta = 0.084, p < .01$) 與越多的 hashtag 使用 ($\beta = 1.096, p < .001$) 分別顯示出與互動呈高度正相關，因此支持 H3 與 H4。Emoji 數量亦具有顯著的正向效果 ($\beta = 0.343, p < .001$)，支持 H5，顯示視覺語言線索能有效提升使用者互動。

\section{討論與管理意涵}
本研究對社群內容策略的文獻做出了實際貢獻，將內容特徵理論落實於 Booking.com 在 X 這個快節奏社群平台上的高度視覺化旅遊情境。藉由驗證理論與收集實證，這項研究替旅遊平台的社群行銷精鍊出明確的管理方針。

H1 慘遭拒絕卻提供了十分重要的管理洞見：在 X 這樣以微網誌為導向的平台上，「長篇大論」只會讓成效適得其反，受眾更傾向快速消化的精簡短語。相較之下，互動提示元件帶來極其優異的表現。適當的問號能引發使用者的強烈參與感並開啟對話；正向情緒吻合旅遊推廣中令人心生嚮往的本質；配置 hashtag 更有助於被捲入當下最熱門的社群觀光討論串中。H5 的成立更進一步證實，emoji 作為視覺語言線索，能豐富品牌內容的表現力並吸引讀者目光，有效帶動互動。

\section{結論}
總結上述，文字特徵本身的設定大幅主導了 Booking.com X 上的貼文參與度。我們給予品牌社群經理的最核心建議是：撰寫言簡意賅、語氣極度正向的文案，並且多半時候以「提問」引領核心訊息來敲開對話大門。有策略地增補 hashtag 以及視覺化的表情符號 (emoji) 更能抓緊讀者目光。

儘管本研究在因變數上等權重地壓縮了留言、按讚與分享行為成一個單一社群互動指標，我們相信強化文本裡的互動特質正是牽動整體受眾行為的關鍵之鑰。未來研究可細分互動屬性或檢視特定旅遊地點是否具備情緒調節效應，但即便是針對當前的每日內容企劃，我們所產出的這套數據觀點也無疑地帶來穩健可靠的文案引導基準。

\section{參考文獻（節錄）}
\begin{enumerate}
    \item De Vries, L., Gensler, S. and Leeflang, P.S.H. (2012). Popularity of brand posts on brand fan pages: An investigation of the effects of social media marketing. \textit{Journal of Interactive Marketing}, 26(2), 83--91.
    \item Gkikas, D. C., Tzafilkou, K., Theodoridis, P. K., Garmpis, A. and Gkikas, M. C. (2022). How do text characteristics impact user engagement in social media posts? Modeling content readability, length, and hashtags number in Facebook. \textit{International Journal of Information Management Data Insights}, 2, 100067.
    \item Ko, E.E., Kim, D. and Kim, G. (2022). Influence of emojis on user engagement in brand-related user-generated content. \textit{Computers in Human Behavior}.
    \item Li, Y. and Xie, Y. (2020). Is a picture worth a thousand words? An empirical study of image content and social media engagement. \textit{Journal of Marketing Research}, 57(1), 1--19.
    \item Mohammad, A.A., Elshaer, I.A., Azazz, A.M., Kooli, C., Algezawy, M. and Fayyad, S. (2024). The influence of social commerce dynamics on sustainable hotel brand image, customer engagement, and booking intentions. \textit{Sustainability}, 16(14), 6050.
    \item Moran, G., Muzellec, L. and Johnson, D. (2019). Message content features and social media engagement: evidence from the media industry. \textit{Journal of Product and Brand Management}, 29(5), 533--545.
    \item Reddy, M.A.S. and Prakash, C. (2024). Sentiment Analysis of Social Media Networking Sites: A Comparative Study on User Engagement and Public Opinion.
    \item van der Harst, J.P. and Angelopoulos, S. (2024). Less is more: Engagement with the content of social media influencers. \textit{Journal of Business Research}, 181, 114746.
\end{enumerate}

\newpage
\appendix
\section{Appendix / 附錄}

% ============================================================
% Table 1 — Descriptive Statistics / 敘述統計
% ============================================================
\begin{table}[htbp]
\centering
\caption{Descriptive Statistics / 敘述統計（N = 766）}
\label{tab:desc}
\small
\begin{tabular}{lrrrrrrrr}
\toprule
Variable / 變數 & Mean & SD & Min & Q1 & Median & Q3 & Max & Skewness \\
\midrule
log(1+Engagement) & 1.942 & 1.572 & 0.000 & 0.693 & 1.609 & 2.996 & 7.814 & 1.001 \\
Text Length (words) & 23.45 & 12.26 & 0.000 & 14.00 & 22.00 & 33.00 & 53.00 & 0.239 \\
Question Mark (0/1) & 0.167 & 0.373 & 0.000 & 0.000 & 0.000 & 0.000 & 1.000 & 1.785 \\
Emotional Valence & 0.789 & 1.528 & $-3.352$ & 0.000 & 0.000 & 2.047 & 3.948 & 0.371 \\
Hashtag Count & 0.145 & 0.489 & 0.000 & 0.000 & 0.000 & 0.000 & 3.000 & 3.817 \\
Emoji Count & 0.398 & 0.884 & 0.000 & 0.000 & 0.000 & 1.000 & 13.000 & 5.261 \\
Picture Count & 0.309 & 0.524 & 0.000 & 0.000 & 0.000 & 1.000 & 4.000 & 1.875 \\
URL Count & 0.919 & 0.719 & 0.000 & 0.000 & 1.000 & 1.000 & 2.000 & 0.121 \\
At-Mention Count & 0.097 & 0.356 & 0.000 & 0.000 & 0.000 & 0.000 & 3.000 & 4.276 \\
\bottomrule
\end{tabular}
\end{table}

\vspace{1em}

% ============================================================
% Table 2 — Pearson Correlation Matrix (lower triangle)
% Note: *** p<.001  ** p<.01  * p<.05
% ============================================================
\begin{table}[htbp]
\centering
\caption{Pearson Correlation Matrix / Pearson 相關矩陣（下三角，N = 766）}
\label{tab:corr}
\small
\setlength{\tabcolsep}{3pt}
\begin{tabular}{lrrrrrrrrr}
\toprule
 & (1) & (2) & (3) & (4) & (5) & (6) & (7) & (8) & (9) \\
\midrule
(1) log(1+Eng.) & 1 & & & & & & & & \\
(2) Text Length & $-0.184^{***}$ & 1 & & & & & & & \\
(3) Question Mk. & $0.262^{***}$ & $-0.191^{***}$ & 1 & & & & & & \\
(4) Emot. Valence & $0.234^{***}$ & $0.150^{***}$ & $0.089^{*}$ & 1 & & & & & \\
(5) Hashtag Ct. & $0.452^{***}$ & $-0.012$ & $0.075^{*}$ & $0.047$ & 1 & & & & \\
(6) Emoji Ct. & $0.411^{***}$ & $-0.044$ & $0.206^{***}$ & $0.257^{***}$ & $0.087^{*}$ & 1 & & & \\
(7) Picture Ct. & $0.451^{***}$ & $-0.027$ & $0.163^{***}$ & $0.312^{***}$ & $0.044$ & $0.423^{***}$ & 1 & & \\
(8) URL Ct. & $0.318^{***}$ & $0.063$ & $-0.018$ & $0.175^{***}$ & $0.193^{***}$ & $0.154^{***}$ & $0.501^{***}$ & 1 & \\
(9) At-Mention Ct. & $0.319^{***}$ & $0.035$ & $0.026$ & $0.056$ & $0.460^{***}$ & $0.127^{***}$ & $0.043$ & $0.102^{**}$ & 1 \\
\bottomrule
\multicolumn{10}{l}{\small $^{***}\,p<.001$,\; $^{**}\,p<.01$,\; $^*\,p<.05$.} \\
\end{tabular}
\end{table}

\vspace{1em}

% ============================================================
% Table 3 — Robust OLS Regression Results
% ============================================================
\begin{table}[htbp]
\centering
\caption{Robust OLS Regression Results predicting Log(Engagement)}
\label{tab:results}
\begin{tabular}{l c}
\toprule
 & \textbf{Coefficient (SE)} \\
\midrule
Intercept & $1.546^{***}$ (0.106) \\
Text Length (H1) & $-0.021^{***}$ (0.003) \\
Question Mark (H2) & $0.479^{***}$ (0.139) \\
Emotional Valence (H3) & $0.084^{**}$ (0.030) \\
Hashtag (H4) & $1.096^{***}$ (0.134) \\
Emoji (H5) & $0.343^{***}$ (0.093) \\
\midrule
Picture & $0.786^{***}$ (0.182) \\
URL & $0.169^{*}$ (0.080) \\
At-Mention & $0.515^{*}$ (0.224) \\
\midrule
Observations & 766 \\
$R^2$ & 0.492 \\
\bottomrule
\multicolumn{2}{l}{\small Notes: Robust (HC) standard errors in parentheses.}\\
\multicolumn{2}{l}{\small $^{***} p<.001$, $^{**} p<.01$, $^* p<.05$.}
\end{tabular}

\vspace{2em}

\caption{穩健 OLS 迴歸結果：預測 $\log(1+\text{Engagement})$}
\label{tab:results_tw}
\begin{tabular}{l c}
\toprule
 & \textbf{係數（穩健標準誤）} \\
\midrule
截距 & $1.546^{***}$ (0.106) \\
文字長度 (H1) & $-0.021^{***}$ (0.003) \\
問號 (H2) & $0.479^{***}$ (0.139) \\
情緒正向程度 (H3) & $0.084^{**}$ (0.030) \\
Hashtag 數量 (H4) & $1.096^{***}$ (0.134) \\
表情符號 (H5) & $0.343^{***}$ (0.093) \\
\midrule
圖片 (Picture) & $0.786^{***}$ (0.182) \\
網址 (URL) & $0.169^{*}$ (0.080) \\
提及 (At-Mention) & $0.515^{*}$ (0.224) \\
\midrule
觀察值數量 & 766 \\
$R^2$ & 0.492 \\
\bottomrule
\multicolumn{2}{l}{\small 註：括號內為穩健標準誤（HC）。}\\
\multicolumn{2}{l}{\small $^{***} p<.001$、$^{**} p<.01$、$^* p<.05$。}
\end{tabular}
\end{table}

\end{document}
